% Transcription — fonds Alexandre Grothendieck, shelfmark 33, batch 1 (pages 1-20).
% Established from the facsimile archives/batches/33/batch-01.pdf (a mirror of
% https://grothendieck.umontpellier.fr/33.pdf), per the skill
% transcribe-grothendieck.
%
% Pass: Opus 5.5 (claude-opus-5-5), 2026-10-03 — first pass, unchecked against the pages by a human.
%
% Eleven of the twenty pages are transcribed: 1, 2, 3, 5, 7, 9, 11, 13, 15,
% 17, 19. The even pages 4, 6, 8, 10, 12, 14, 16, 18 and page 20 are typed
% versos (a Bourbaki-style draft on root systems and Dynkin diagrams, pages
% numbered 179-194, « n° 391 ») carrying nothing in his hand, and are absent.
%
% The batch is one continuous programme, in ink on the backs of those
% typescript sheets, for a course or seminar on étale cohomology: « Nous
% prendrons comme prétexte la théorie des fonctions L … ». Numbered items 1° to
% 14° run through the batch: étale site, functoriality, fibre functors, strict
% localisation, the proper base change theorem, comparison with classical
% cohomology, specialisation, cohomological dimension, supports, purity,
% constructible sheaves, and the R^q_! f_* of a compactifiable morphism, whose
% list of properties is still open at the foot of page 19. Page 2 is a sheet
% of pencil diagrams (base-change squares) in the same hand.
%
% The hand is fast and heavily abbreviated; formulas come through, the
% connecting prose often does not. His underlinings are set with \emph.
% Page 13's item « 13) Les R^i_! f(F) » and its continuation at the head of
% page 15 are cancelled by large crossing strokes; the cancellation is
% recorded once, and the text is rewritten as 14) on page 17.

\documentclass[11pt,a4paper]{article}
\input{../preamble/grothendieck}

\folder{33}
\batch{1}
\pages{1}{20}
\dating{1965-[vers 1971]}
\watermark{Édition de démonstration}
\foldertitle{SGA 1965. MS [Manuscrit] brouillon : notes manuscrites (s.d.).}

\begin{document}

\section{Programme : cohomologie étale et fonctions $L$}

\page{1}
Nous prendrons comme prétexte la théorie des fonctions $L$ pour développer plus
avant le formalisme de la cohomologie étale. Trois parties

a) Continuer l'étude \add{de la coh. étale} avec des coefficients de torsion,
en \struck{développant} en particulier, dans ce cadre, \struck{\ill{}}
\uncertain{avant} le formalisme de dualité, \struck{donnant} donnant des
formules de Künneth et de Lefschetz.

b) Étendre le formalisme : aux coefficients \emph{$\ell$-adiques}.

c) Appliquer les résultats au formalisme des fonctions $L$.

1°) \emph{Site étale} de $X$,

$X_{\mathrm{ét}}$ : catégorie, familles couvrantes, faisceaux, faisceaux
abéliens, cohomologie

2°) \emph{Exemples} $X = \operatorname{Spec}(k)$,
$X_{\mathrm{ét}} \simeq \widetilde{X}_{\mathrm{ét}} \simeq \mathrm{Ens}(\pi)$,
\qquad $\pi = \operatorname{Gal}(\bar{k}/k)$
\note{le premier $\simeq$ est tel quel sur la page, entre le site et le topos.}

3°) $f : X \to Y$ définit $X_{\mathrm{ét}} \to Y_{\mathrm{ét}}$
\[
  f^{*} : Y_{\mathrm{ét}} \to X_{\mathrm{ét}}, \qquad
  \widetilde{Y}_{\mathrm{ét}} \to \widetilde{X}_{\mathrm{ét}}, \qquad
  f_{*} : \widetilde{X}_{\mathrm{ét}} \to \widetilde{Y}_{\mathrm{ét}}
\]
\marginal{pour $f^{*}$ : commute aux \emph{lim.\ ind.} quelc.\ et aux
\emph{lim.\ proj.} finies ; pour $f_{*}$ : commute aux lim.\ proj.\ quelc.}

\page{2}
\drawing{feuille de croquis au crayon. En haut à droite « McL » et
« $\times M'U$ », une lettre « A ». À gauche, un morphisme
$X \to k$ et $k \to Y$ en équerre, puis $X' \to X$ avec une flèche oblique
$X' \to k$ marquée \uncertain{« $\phi$ »}. En bas, deux esquisses
$Y \to X$ au-dessus de $k$ et $Y' \to X'$ au-dessus de $k'$, une flèche
$k \to k'$ barrée d'un petit cercle hachuré ; à droite un carré
$X \leftarrow X$, $K$, $k'$, $k \leftarrow K$, un $X$ entouré.}

Au milieu à droite, un carré entouré :

\begin{tikzcd}
Y \arrow[d, "f"'] & Y' \arrow[l, "h"'] \arrow[d, "{f'}"] \\
X & X' \arrow[l, "g"]
\end{tikzcd}

\note{à côté du carré : « $V'$ » près de $Y'$, « $g(U')$ » sous $X$, et
« $U' \neq \emptyset$ » près de $X'$.}

\page{3}
transitivité …

\emph{Suite spectrale de Leray} \add{$R^i f_{*}$}

4°) Foncteurs fibres

Soit $\xi$ un pt géométrique de $X$
\[
  \xi \to X, \quad \text{d'où foncteur fibre} \quad
  u^{*} : \widetilde{X}_{\mathrm{ét}} \to \widetilde{\xi}_{\mathrm{ét}} \simeq (\mathrm{Ens})
\]
La famille des foncteurs fibres est \emph{conservative} : $u : F \to G$ est un
isom ssi les $u_\xi : F_\xi \to G_\xi$ le sont pour tt $\xi$
\uncertain{$\in X$}, d'où les cor.\ habituels.
\[
  F_\xi = \varinjlim_{\substack{X' \text{ étale sur } X\\ \text{et pointé par } \xi}} F(X')
\]

5°) \emph{Localisation stricte}
\[
  \mathcal{O}_{X,\xi} = \varinjlim_{\substack{X' \text{ étale sur } X\\ \text{et } \xi\text{-pointé}}} \Gamma(X', \mathcal{O}_{X'})
\]
\note{à gauche un petit diagramme : $\xi \to X'$, $X' \to X$, et $\xi \to X$ ;
un symbole biffé précède le signe $=$.}
anneau des stricts localisés en $\xi$ [anneau \struck{\ill{}}
\add{strict.\ local} i.e.\ hensélien et à corps rés.\ sép.\ clos]
\[
  \overline{X}{}^{\xi} = \operatorname{Spec} \mathcal{O}_{X,\xi}, \qquad
  F_\xi = \Gamma(\overline{X}, F_{\overline{X}})
\]

6°) Application aux fibres des $R^q f_{*}(F)$, pour $f$ qu.-cpt.\ et qu.-sép.
\marginal{NB 5° et 6° n'ont pas vraiment d'analogue complet dans le cas des
espaces top.\ ordinaires …}

7°) \emph{Théorème fondamental} Soit $f : X \to Y$ morphisme propre, $F$
faisceau \emph{de torsion} sur $X$, alors la formation des $R^q f_{*}(F)$

\page{5}
commute à tt changement de base $Y' \to Y$.

\emph{NB} Vrai sans hyp.\ de propreté si $Y' \to Y$ est \emph{lisse}, $f$
qu.-cpt.\ qu.-sép., et $F$ premier aux car.\ rés.\ de $Y$.

\struck{8°) Ex. Morphismes finis : $R^q f_{*}(F) = 0$ pour $q \geq 1$}
\note{ligne biffée, lecture approximative.}

Ex. $R^q f_{*}(F)_{\bar y} \simeq H^q(X_{\bar y}, F_{\bar y})$

Ex. $X$ propre sur $k$ sép.\ clos, $K$ ext.\ sép.\ close de $k$, d'où
$X' = X_{K}$, $F' = F_{X'}$, et $H^i(X, F) \simeq H^i(X', F')$

\emph{NB} Vrai sans hyp.\ de propreté si $F$ de torsion premier à
\struck{\ill{}} car $k$.

8°) \emph{Comparaison avec la cohomologie classique}

$X$ préschéma loc.\ de type fini sur $\mathbb{C}$, alors pour tt $F$ faisceau
de torsion sur $X$,
\[
  H^i(X, F) \to H^i(X^{\mathrm{an}}, F^{\mathrm{an}})
\]
est un \emph{isomorphisme}. Résultat analogue pour les $R^i f_{*}$ pour un
morphisme $f : X \to Y$, si $f$ est \emph{de type fini}.

\note{le paragraphe suivant est biffé de quatre grands traits obliques ; il
est lu ici en partie.}
\struck{Soit $X$ un schéma propre sur un anneau $A = \mathcal{O}_{Y,y}$ \ill{}
$\mathbb{C}$, \ill{} $Y$ \ill{}. Alors $X$ « provient » d'un \ill{}
algébrique sur $\mathbb{C}$ \ill{} $X'$ propre \ill{} au-dessus d'un voisinage
\ill{} $Y'$ de $Y$. \ill{} faisceau de torsion constr.\ $F'$ sur $X'$
« induisant » un faisceau constr.\ $F$ sur $X$.}

\marginal{Utiliser de plus \uncertain{résolution des singularités}
\struck{\ill{}} et l'hyp.\ de propreté.}
\marginal{\emph{GAGA} plus complet : OK. \struck{\ill{}} Ext.\ pour faisceaux
de torsion. Équivalence de catégories des faisceaux constructibles algébriques
et faisceaux constructibles analytiques \uncertain{se réduisant} de torsion.
\ill{} relatives …}

\page{7}
\struck{On peut certainement \ill{} que
$H^i(X, F) \simeq \varinjlim_{U \ni y} H^i(f^{-1}(U), F)$}
\add{Il y a aussi des \ill{} relatifs \uncertain{dans le type} :}

\emph{Remarque} Soit $f : X \to Y$ un « schéma relatif » de type fini sur un
espace analytique $Y$, d'où $f^{\mathrm{an}} : X^{\mathrm{an}} \to Y$. Tout
faisceau étale \add{de torsion} $F$ constructible sur $X$ définit un faisceau
analytique $F^{\mathrm{an}}$ sur $X^{\mathrm{an}}$, d'où les hom.
\[
  H^q(X, F) \to H^q(X^{\mathrm{an}}, F^{\mathrm{an}}), \qquad
  R^q f_{*}(F)^{\mathrm{an}} \to R^q f^{\mathrm{an}}_{*}(F^{\mathrm{an}})
\]
À prouver que ce sont des \emph{isomorphismes}.
\marginal{Devrait aussi se généraliser au contexte rigide-analytique …}

\emph{Var.} Passant à la limite sur les voisinages d'un pt $y$ de $Y$, on
trouverait alors
\[
  H^i(X \times_{Y} \mathcal{O}_{Y,y}, F \otimes \mathcal{O}_{Y,y})
  \simeq \varinjlim_{U \ni y} H^i((f^{\mathrm{an}})^{-1}(U), F^{\mathrm{an}})
\]
\note{sous le premier membre, une flèche vers l'annotation « schéma ordinaire
de t.f.\ sur $\mathcal{O}_{Y,y}$ ».}

Par exemple, prenons $X = Y - Z$, $Z$ fermé analytique, en vérifiant
\uncertain{par là} la \uncertain{classe} des $R^q f_{*}(F)$ ($F$ faisceau de
torsion, \uncertain{loc.\ cst}, sur $X - Z$) \struck{\ill{}} provient d'un
faisceau dans le cadre de la \emph{géom.\ alg.}

\page{9}
9°) \emph{Th.\ de spécialisation} Voici le type d'une forme précise :
$f : X \to Y$ morphisme \emph{propre} et \emph{lisse}, $F$ faisceau de
torsion loc.\ libre sur $X$, premier aux car.\ résiduelles de $X$, alors les
$R^q f_{*}(F)$ [dont les fibres sont les coh.\ des fibres …] sont loc.\ libres.

10°) \emph{Dim.\ cohomologique.}

$X$ de type fini sur $k$ sép.\ clos, $n = \dim X$, alors
$H^i(X, F) = 0$ pour $i$ \struck{$\geq$} $> 2n$,
($F$ faisceau de torsion sur $X$).
\note{un $\geq$ biffé est corrigé en $>$.}

Si $X$ est \emph{affine}, on a même $H^i(X, F) = 0$ pour $i > n$ (th.\ du
type Lefschetz).

La dimension « topologique » de $X$ \uncertain{devant} donc être en général le
double de la dimension de Krull — il y a un abaissement dans le cas affine, qui
peut donner des idées fausses.

\page{11}
(cf Artin, Woodshole, par changement de base)

11) \emph{Supports} \struck{Désignons par $U$ ouvert $U$ de $X$,}

Faisceau final : $X/X$ lui-même, $X(X') = \{e\}$ ; les sous-faisceaux du
faisceau final \struck{sont} sont les ouverts de $X$.

\struck{Support d'un f}
\struck{Si $F$ est un fai}
La catégorie des faisceaux $F$ sur $X$ qui sur $U$ sont le faisceau final est
équivalente à la catégorie des faisceaux sur $Y = X - U$, par le foncteur image
\add{directe}. D'où : faisceaux abéliens sur $X$ nuls sur $U$ $=$ faisceaux
abéliens sur $Y$.

Support d'un faisceau abélien
\struck{$\operatorname{Supp}(F) \subset \Gamma(F)$ fermé.}
Support d'une section d'un faisceau abélien
\[
  \Gamma_Y(F) \subset \Gamma(F), \qquad
  \underline{\Gamma}_Y(F) \subset F, \qquad
  \underline{\Gamma}_Y(F)(X') = \Gamma_{Y'}(X', F')
\]

Foncteurs dérivés
\[
  H^i_Y(X, F) \simeq \operatorname{Ext}^i(X ; \mathbb{Z}_Y, F), \qquad
  \underline{H}^i_Y(F) \simeq \underline{\operatorname{Ext}}^i(\mathbb{Z}_Y, F)
  = R^{i-1} u_{*}(F|U) \quad \text{si } i \geq 2
\]
\marginal{si $i = 0$ ou $1$, noyau et conoyau de $F \to u_{*}(F|U)$}
Suite exacte, spectrale …
\marginal{cf SGA 1962, notes Hartshorne 1961}

12) \emph{Pureté} $Y \subset X$ régulière, codim${}_X(Y) = p$ en chaque pt
de $Y$

\begin{tikzcd}
Y \arrow[rr, hook] \arrow[dr, "\text{lisse}"'] & & X \arrow[dl, "\text{lisse}"] \\
& k &
\end{tikzcd}

\[
  \underline{H}^i_Y((\mathbb{Z}/n\mathbb{Z})_X) =
  \begin{cases}
    0 & \text{si } i \neq 2p \\
    \simeq \mu_n^{\otimes(-p)} & \text{si } i = 2p
  \end{cases}
  \qquad n \text{ premier aux car.\ résiduelles de } Y
\]
\note{l'exposant se lit $\otimes(-p)$ sur la page.}
\marginal{Donner explication de \struck{lisse} la classe fondamentale}

\page{13}
\emph{NB}
\drawing{un cercle traversé d'un diamètre horizontal, un point au centre.}
\[
  \mu_n = \operatorname{Ker}(\mathbb{G}_m \xrightarrow{\;n\;} \mathbb{G}_m),
  \qquad
  \mu_n(X') = {}_n\Gamma(X', \mathcal{O}_{X'})^{*}
\]
\note{un « Ker » est biffé devant ${}_n\Gamma$.}
\uncertain{gr.} fini plat sur $X$.

Si $n$ premier aux car.\ de $X$, $\mu_{n,X}$ est \emph{étale} sur $X$, loc.\
isom.\ (pour top.\ étale) à $\mathbb{Z}/n\mathbb{Z}$. Il \add{ne} faut surtout
pas les confondre. Les $\mu_{n,X}^{\otimes d}$ ($d \in \mathbb{Z}$) sont donc
définis, et sont \add{loc.} isomorphes à $(\mathbb{Z}/n\mathbb{Z})_X$ [le
choix d'une section génératrice de $\mu_n$ définit un isom
$\mu_{n,X}^{\otimes d} \simeq (\mathbb{Z}/n\mathbb{Z})_X$, multiplié par
$\lambda^d$ si on multiplie la section par
$\lambda \in (\mathbb{Z}/n\mathbb{Z})^{*}$ \ill{} $\lambda^d$ …].
\emph{Dém.} pureté immédiate par th.\ de changement de base …

\emph{Pb} [Pureté « absolue »] avec hypothèse de lissité remplacée par hyp.\
régularité. OK \uncertain{pour toutes puissances} [\uncertain{pour} \ill{}
\ill{}] \ill{} dépend de la résolution des singularités (Artin).

Il y a aussi un th.\ de pureté relatif sur une base quelconque …

\note{la suite de la page et le haut de la p.\ 15 sont barrés de grands traits
croisés ; l'article est repris en 14) p.\ 17.}
\struck{13) B) \emph{Les} $R^i_{!} f(F)$. Morphismes compactifiables
$f : X \to Y$ \ill{} \ill{} Ext.\ aux \ill{} \ill{} croisements normaux …}

\page{15}
\struck{\ill{} ouvert, $\bar{f}$ propre \ill{} \ill{} quasi-projectif \ill{}
st.\ \ill{} \ill{}.}

13) \emph{Faisceaux constructibles}

Faisceaux loc.\ \struck{\ill{}} constants

Faisceaux constructibles : en général, il \uncertain{a}
loc.\ constant : fibres \struck{finies} \add{de t.f.}
sur des ens.\ loc.\ fermés qui partitionnent $U$. Dans le cas noethérien,
\uncertain{critère} par adhérences \ill{} \struck{\ill{}}.

\emph{Th.} Si $f : X \to Y$ propre, $Y$ loc.\ \emph{noethérien}, alors pour tt
faisceau \emph{de torsion} \emph{constructible} $F$ sur $X$, les
$R^i f_{*}(F)$ sur $Y$ sont constr.

\emph{Cor.} $X$ propre sur un \uncertain{corps sép.\ clos}, \ill{} $F$ de
torsion constructible sur $X$, les $H^i(X, F)$ sont finis.

\emph{Remarques} \uncertain{Questions} de la généralisation

\page{17}
: au $f$ de type fini au lieu de $f$ propre. Se ramène au cas d'une
\emph{immersion}. Se ramène \uncertain{par} la \uncertain{résolution} des
singularités \add{+ pureté} \add{ou plutôt « bons modèles »},
\uncertain{lorsque} \ill{}. Cf.\ là-dessus \uncertain{la} \ill{}
\uncertain{bons} \uncertain{modèles} \ill{} des singularités \struck{\ill{}}
dans ce cadre …

14) \emph{Les} $R^i_{!} f(F)$.

Soit $f : X \to Y$ morphisme compactifiable, i.e.\ qui se factorise en
\[
  X \xrightarrow{\;i\;} \overline{X} \xrightarrow{\;\bar{f}\;} Y
\]
$i$ immersion ouverte, $\bar{f}$ propre (cas où $f$ qu.-proj.\ et $Y$
\uncertain{qu.-cpt}).

Pour tt faisceau de torsion $F$ sur $X$, on considère les $i_{!}(F)$ [ext.\ de
$F$ par $0$], puis les $R^q \bar{f}_{*}(i_{!}(F))$, foncteurs coh.\ en $F$.
Je \uncertain{vérifie} \uncertain{qu'ils} ne dépendent \ill{} de façon
\ill{} \ill{} \ill{} fonctions \uncertain{cohomologiques}
\ill{} \ill{} \ill{} foncteur coh.

\page{19}
\uncertain{à} \ill{} \uncertain{près}, \ill{} \emph{indépendant} du choix de
la factorisation.
\emph{Notation} $R^q_{!} f_{*}(F)$.
\[
  R^q_{!} f_{*}(F) = (R^q \bar{f}_{*})(i_{!}(F))
\]
\note{lecture du membre de droite douteuse ; suit sur la même ligne
« \ill{} : \ill{} propre sur $Y$ ».}

\struck{Propriétés} Propriétés fondamentales, \uncertain{identiques} à celles
des $R^q f_{*}(F)$, \emph{aussi bien} que si $f$ propre :

(i') \emph{Comparaison avec cas classiques sur $\mathbb{C}$}

(ii) \emph{Th.\ de changement de base} \struck{\ill{}}, \emph{dito.}

(iii) \add{\emph{Th.\ de finitude}} \struck{Suite spectrale de
Leray}

(iv) \emph{Suite spectrale de Leray}

(v) Si $f$ quasi-fini, \add{\uncertain{alors}} \ill{} les
$R^i_{!} f_{*}(F) = 0$ \uncertain{pour $i > 0$}
\marginal{phénomènes spéciaux}

(vi) \struck{D'autre part} Si $f$ \emph{étale}, on a un hom.\
\uncertain{can.}
\[
  \operatorname{Tr}_f(F) : f_{!}(F) \to F
\]
[d'où les covariants, \uncertain{par} $f : X' \to X$ morphisme
\ill{}, \ill{} \ill{} \ill{} \ill{} \uncertain{compactifiable} sur $Y$
\[
  R^q_{!} g'_{*}(F') \to R^q_{!} g_{*}(f_{!}(F'))
\]
on en \uncertain{tire}
\[
  R^q_{!} g'_{*}(f^{*}(F)) \to R^q_{!} g_{*}(F) \;]
\]

(vi) \emph{Suite exacte en cohomologie relative}
\note{deux articles portent le numéro (vi) sur la page. La liste se poursuit
au-delà de la p.\ 19, dans le lot suivant.}

\end{document}
